% ============================================================
\part{Vorbereitung auf ROS2}
% ============================================================

\chapter{Wir wollen verstehen, warum ROS2 überhaupt existiert}

\kapitelziel{
Du verstehst die Motivation hinter ROS2 und kannst erklären,
welche Probleme es löst -- und welche es nicht löst.
}
\kapitelstruktur

\section{Situation}
Du hast jetzt einen funktionierenden Roboter: Python-Code für
Sensoren und Motoren, Shell-Skripte für den Start, Docker für
die Umgebung. Warum braucht man noch ROS2?

\begin{aufgabeBox}
Verstehe, welche Probleme bei größeren Robotikprojekten entstehen
und warum ROS2 diese systematisch löst.
\end{aufgabeBox}

\begin{problemBox}
Ohne Framework muss jede Verbindung manuell programmiert werden.
Bei 7~Sensoren und 5~Empfängern sind das 35~individuelle Verbindungen.
Zeitstempel-Synchronisation, Visualisierung und Multi-Rechner-Betrieb
erfordern jedes Mal eigenen Code.
\end{problemBox}

\begin{wissenBox}
ROS2 bietet: Topics (Publish-Subscribe), Services (Anfrage-Antwort),
Actions (Langläufer mit Feedback), Parameter, tf2
(Koordinatentransformationen), rviz2 (Visualisierung),
rosbag (Aufzeichnung), Gazebo-Integration (Simulation).
\end{wissenBox}

\begin{lstlisting}[style=text]
Ohne ROS2:                        Mit ROS2:
sensor.py -> navigation.py        /scan Topic
sensor.py -> anzeige.py      ==>  /scan Topic (alle subscriben)
sensor.py -> karte.py             /scan Topic
Jede Verbindung: eigener Code     Eine Nachricht, viele Empfaenger
\end{lstlisting}

\begin{merksatzBox}
ROS2 ist kein Betriebssystem und keine Programmiersprache.
Es ist ein Middleware-Framework: Es kümmert sich um die
Kommunikation zwischen Programmen und stellt Werkzeuge bereit.
Den Algorithmus musst du selbst schreiben.
\end{merksatzBox}

\begin{robotikBox}
ROS2 nutzt DDS (Data Distribution Service) als Transport-Layer --
einen ISO-Standard, der auch in der Luft- und Raumfahrt sowie
Medizintechnik eingesetzt wird. Das macht ROS2 deterministisch,
skalierbar und netzwerkfähig über mehrere Rechner hinweg.
\end{robotikBox}

% ── Kapitel 26 ────────────────────────────────────────────────
\chapter{Wir bauen ein Mini-ROS im Kopf}

\kapitelziel{
Du kannst das Zusammenspiel von Nodes, Topics, Messages und
Services anhand eines konkreten Roboterszenarios erklären.
}
\kapitelstruktur

\section{Situation}
Bevor du ein echtes ROS2-System mit Dutzenden Nodes siehst,
baust du ein vereinfachtes Modell im Kopf.

\begin{aufgabeBox}
Beschreibe den ROS2-Stack eines einfachen mobilen Roboters
mit Nodes, Topics, Messages und einem Service.
\end{aufgabeBox}

\begin{lstlisting}[style=text]
Nodes und Topics:
+------------------+  Topic: /scan (Float32)  +------------------+
| ultraschall_node | -----------------------> | hindernis_node   |
+------------------+                          +------------------+
                                                       |
                                             Topic: /cmd_vel (Twist)
                                                       |
+------------------+  Topic: /cmd_vel         +------------------+
|  dashboard_node  | <----------------------- |   motor_node     |
+------------------+                          +------------------+
       |
  Service: /get_status [std_srvs/Trigger]
\end{lstlisting}

\begin{lstlisting}[style=text]
Message-Typen:
Float32 (std_msgs):     Twist (geometry_msgs):
  float32 data            Vector3 linear   -> linear.x: 0.3 m/s
  z.B.: data: 0.42        Vector3 angular  -> angular.z: 0.0 rad/s
\end{lstlisting}

\begin{merksatzBox}
\textbf{Node} = ein Programm (Python-Klasse, erbt von \texttt{Node}).\\
\textbf{Topic} = benannter Kanal, z.\,B. \texttt{/scan}.\\
\textbf{Message} = Datenstruktur, die über Topic fließt.\\
\textbf{Publisher} = schreibt auf Topic.\\
\textbf{Subscriber} = liest von Topic.\\
\textbf{Service} = synchrone Anfrage-Antwort.
\end{merksatzBox}

\begin{robotikBox}
Das Pub-Sub-Muster aus Kapitel 22 (eigener MessageBus) ist
konzeptionell identisch mit ROS2-Topics. ROS2 fügt hinzu:
automatische Discovery, typsichere Messages und netzweite
Kommunikation -- das Prinzip bleibt dasselbe.
\end{robotikBox}

% ── Kapitel 27 ────────────────────────────────────────────────
\chapter{Wir richten die Robotik-Werkstatt vollständig ein}

\kapitelziel{
Du hast eine vollständige, dokumentierte Entwicklungsumgebung
für ROS2 auf Basis von Docker -- bereit für Band 2.
}
\kapitelstruktur

\section{Situation}
Alle Werkzeuge sind vorhanden. Jetzt fügen wir sie zur
vollständigen Robotik-Werkstatt zusammen.

\begin{aufgabeBox}
Richte die komplette Entwicklungsumgebung ein:
ROS2 via Docker, VS Code Remote, Git, udev für Arduino.
\end{aufgabeBox}

\begin{lstlisting}[style=bash]
# 1. Docker einrichten
sudo apt install docker.io docker-compose
sudo usermod -aG docker $USER && newgrp docker

# 2. Projektstruktur
mkdir -p ~/robotik/{config,src,scripts,logs}
cd ~/robotik && git init

# 3. docker-compose.yml anlegen (nano oder VS Code)
# Inhalt: ros:jazzy Image, Volumes fuer src/ und logs/

# 4. ROS2-Container testen
docker-compose run ros2 bash -c \
    "source /opt/ros/jazzy/setup.bash && ros2 --version"

# 5. udev-Regel fuer Arduino
echo 'SUBSYSTEM=="tty",ATTRS{idVendor}=="2341",SYMLINK+="arduino",MODE="0666"' \
    | sudo tee /etc/udev/rules.d/99-arduino.rules
sudo udevadm control --reload-rules

# 6. SSH-Key fuer VS Code Remote
ssh-keygen -t ed25519 -f ~/.ssh/robotik_key
\end{lstlisting}

\begin{lstlisting}[style=yaml]
# docker-compose.yml fuer die Werkstatt
version: '3.8'
services:
  ros2:
    image: ros:jazzy
    volumes:
      - ./src:/robotik/src
      - ./config:/robotik/config
      - ./logs:/root/.ros/log
    working_dir: /robotik
    stdin_open: true
    tty: true
    network_mode: host
\end{lstlisting}

\begin{loesungBox}
\textbf{Täglicher Workflow:}
\begin{lstlisting}[style=bash]
cd ~/robotik
git pull                            # Aktuellen Stand holen
docker-compose run ros2 bash        # ROS2-Umgebung betreten
source /opt/ros/jazzy/setup.bash    # ROS2 aktivieren
# ... entwickeln und testen ...
exit
git add . && git commit -m "..."    # Versionieren
git push                            # Sichern und teilen
\end{lstlisting}
\end{loesungBox}

\begin{merksatzBox}
Die Werkstatt ist fertig, wenn: ROS2 im Container läuft,
Code versioniert in Git liegt, Hardware per udev stabil
erreichbar ist und Remote-Entwicklung über SSH möglich ist.
Erst dann beginnt Band 2.
\end{merksatzBox}

% ── Kapitel 28 ────────────────────────────────────────────────
\chapter{Ausblick: Was passiert in Band 2 und Band 3?}

\kapitelziel{
Du weißt, was dich in den nächsten Bänden erwartet, und
erkennst, wie Band 1 das Fundament für alles Folgende legt.
}

\section{Was du jetzt kannst}

Du hast in Band 1 gelernt:
\begin{itemize}
  \item Linux navigieren, Dateien verwalten, Prozesse steuern
  \item Konfigurationsdateien bearbeiten, Skripte schreiben
  \item Projekte strukturieren und mit VS Code remote entwickeln
  \item Git für Versionsverwaltung und GitHub für Zusammenarbeit
  \item Docker für reproduzierbare Umgebungen
  \item IPC-Grundlagen und das Publish-Subscribe-Muster
  \item USB-Hardware einbinden und mit Arduino kommunizieren
  \item ROS2-Konzepte: Nodes, Topics, Messages, Services
\end{itemize}

Das ist kein Vorwissen -- das ist das Fundament.

\section{Band 2: Python für Robotik}

In Band 2 programmierst du deinen MobileRobot-1 in Python:
Sensorauswertung, Bewegungssteuerung, Hinderniserkennung, Navigation.
Und am Ende: dein erster eigener ROS2-Node.

\begin{lstlisting}[style=python]
# Vorschau: Band 2 -- erster ROS2-Node
class MobilRoboter(Node):
    def __init__(self):
        super().__init__('mobile_robot')
        self.sub = self.create_subscription(
            Float32, '/scan', self.scan_callback, 10)
        self.pub = self.create_publisher(
            Twist, '/cmd_vel', 10)

    def scan_callback(self, msg):
        cmd = Twist()
        if msg.data < 0.25:
            cmd.linear.x = 0.0   # STOPP
        else:
            cmd.linear.x = 0.3   # vorwaerts
        self.pub.publish(cmd)
\end{lstlisting}

\section{Band 3: ROS2 vertieft}

Band 3 zeigt den vollständigen ROS2-Stack:
SLAM (Karte erstellen), Nav2 (autonom navigieren),
Sensorfusion, Launch-Dateien und Deployment auf dem echten Roboter.

\begin{merksatzBox}
Die Robotik-Werkstatt ist keine Vorstufe -- sie ist der Boden,
auf dem alles Folgende steht. Wer Linux, Git und Docker beherrscht,
bewegt sich in ROS2 wie zu Hause.
\end{merksatzBox}

\begin{robotikBox}
\textbf{Dein MobileRobot-1 nach Band 3:}
\begin{itemize}
  \item Fährt autonom und weicht Hindernissen aus
  \item Erstellt eine Karte seiner Umgebung (SLAM)
  \item Navigiert zu einem vorgegebenen Zielpunkt (Nav2)
  \item Wird von ROS2 auf Raspberry Pi gesteuert
  \item Startet nach dem Einschalten automatisch
\end{itemize}
Das ist das Ziel. Band 1 war der erste Schritt.
\end{robotikBox}
